% !TEX root = CoeneDylanBP.tex

\section{ToF-schatting}
\label{sec:tof}

Naast de horizontale verplaatsing is de vluchttijd (ToF) een tweede kerncomponent van de eindscore. Het doel van deze sectie is om de ToF te bepalen op basis van de gedetecteerde landingsmomenten en de framerate van de inputvideo, zonder dat er extra sensoren aan te pas komen.

Omdat de video's vanuit een zijdelings perspectief zijn gefilmd, is een directe schatting van de absolute vluchthoogte ($z$-co{\"o}rdinaat) niet haalbaar zonder dieptesensor. In plaats daarvan wordt de ToF bepaald via \textit{contactdetectie}: de landingsframes gedetecteerd in Sectie~\ref{sec:landingsmoment} fungeren als ankerpunten voor de tijdsmeting.

De ToF-berekening verloopt als volgt: voor elke sprong in een reeks wordt het frame-verschil bepaald tussen het landingsframe van de vorige sprong en het landingsframe van de volgende sprong. Dit verschil, gedeeld door de framerate van de video (frames per seconde), geeft de duur van één vluchttijdsinterval. De totale ToF-score is de som van alle vluchttijdsintervallen over de volledige reeks van tien sprongen.

De nauwkeurigheid van deze methode is inherent begrensd door de framerate van de video. Elke landing kan slechts met een precisie van $\tfrac{1}{\text{fps}}$ seconden worden bepaald. Bij een framerate van 25 fps bedraagt deze onzekerheid 40\,ms per landing, wat over een volledige reeks kan oplopen tot een cumulatieve fout van honderden milliseconden. Dit wordt verder besproken in de conclusie.
